\section{Giới thiệu}
\label{sec:introduction}

Sự cố tràn dầu trên biển là một trong những thảm họa sinh thái và môi trường nghiêm trọng nhất, đe dọa trực tiếp tới đa dạng sinh học biển, các rạn san hô, ngành đánh bắt thủy hải sản và kinh tế du lịch ven biển~\cite{fingas2014review}. Điển hình là các sự cố tràn dầu quy mô lớn trong lịch sử như sự cố nổ giàn khoan Deepwater Horizon năm 2010 tại Vịnh Mexico làm rò rỉ hàng trăm triệu lít dầu thô, hay các vụ rò rỉ từ cơ sở hạ tầng dầu khí Al-Khafji năm 2017 tại Vịnh Ba Tư~\cite{zhu2022oil}. Đối với các quốc gia có bờ biển dài như Việt Nam, việc phát hiện sớm và khoanh vùng chính xác các vệt dầu loang ngoài khơi đóng vai trò quyết định trong việc triển khai các biện pháp ứng phó khẩn cấp và bảo vệ hệ sinh thái biển.

Trong công tác giám sát hàng hải diện rộng, Radar khẩu độ tổng hợp (Synthetic Aperture Radar - SAR) đặt trên vệ tinh (như Sentinel-1 của Cơ quan Vũ trụ Châu Âu ESA hay ALOS PALSAR của Cơ quan Thám hiểm Hàng không Vũ trụ Nhật Bản JAXA) là công cụ viễn thám ưu việt nhất hiện nay nhờ khả năng chụp ảnh chủ động xuyên mây, hoạt động liên tục cả ngày lẫn đêm và không phụ thuộc vào điều kiện thời tiết hay ánh sáng mặt trời~\cite{brekke2005oil,topouzelis2008oil}. Cơ chế vật lý cơ bản của việc phát hiện vệt dầu trên ảnh SAR dựa trên sự tương tác giữa sóng điện từ và mặt biển: lớp màng dầu mỏng nổi trên bề mặt làm tăng độ nhớt động học, dập tắt các sóng mao dẫn và sóng trọng lực ngắn (sóng Bragg) do gió tạo ra. Hiện tượng này làm phẳng mặt biển cục bộ, khiến bức xạ radar bị phản xạ gương ra khỏi hướng thu của anten vệ tinh, dẫn đến cường độ tán xạ ngược (backscatter) bị suy giảm mạnh. Do đó, các vệt dầu thường xuất hiện dưới dạng các vùng tối (dark formations) nổi bật trên nền biển xung quanh~\cite{yang2025oil}.

Tuy nhiên, việc tự động hóa quá trình phát hiện vệt dầu loang từ ảnh SAR đối mặt với ba thách thức kỹ thuật lớn:
\begin{enumerate}
    \item \textbf{Nhiễu đốm nhân (Multiplicative Speckle Noise):} Do bản chất tạo ảnh giao thoa kết hợp của sóng radar, ảnh SAR luôn bị bao phủ bởi nhiễu đốm hạt có độ lệch chuẩn tỷ lệ thuận với cường độ tín hiệu, làm nhòe biên và gây dao động mức xám dữ dội, khiến các thuật toán phân ngưỡng thông thường dễ bị vỡ hạt và sinh hàng loạt đốm vụn giả mạo~\cite{lee1980digital}.
    \item \textbf{Các vùng tối giả mạo trên biển (Dark Formations / False Positives):} Các hiện tượng tự nhiên như vùng biển lặng gió (vận tốc gió $< 2-3\text{ m/s}$ không đủ kích hoạt sóng mao dẫn), màng hữu cơ sinh học do tảo biển tiết ra, dòng chảy nội hay vùng đổ bóng mưa cũng làm suy giảm tán xạ ngược, tạo ra các vùng tối có độ tương phản mức xám tương tự vệt dầu loang~\cite{brekke2005oil,solberg2007oil}.
    \item \textbf{Sự biến động quy mô diện tích và hình học:} Vết dầu trên biển có thể tồn tại dưới dạng dải hẹp phân mảnh chỉ vài pixel cho đến những mảng loang khổng lồ bao trọn toàn bộ khung ảnh viễn thám, gây khó khăn cho các phương pháp phân ngưỡng đơn lẻ.
\end{enumerate}

Để giải quyết các thách thức trên với một kiến trúc nhẹ, trong suốt và có thể diễn giải về mặt vật lý, nghiên cứu này xây dựng và đánh giá thực nghiệm một quy trình xử lý ảnh tuần tự kết hợp lọc đốm thống kê, phân đoạn ứng viên và học máy phân loại theo vùng (Sequential Statistical Filtering, Candidate Thresholding, and Region Classification Pipeline). Các đóng góp chính của nghiên cứu gồm:
\begin{itemize}
    \item Thiết kế và hiện thực hóa một quy trình tuần tự 2 giai đoạn (6 khối chức năng) hoàn chỉnh, độc lập với nhãn tham chiếu ground truth trong suy luận, kết hợp lọc thống kê cục bộ, phân ngưỡng Otsu, lọc diện tích và học máy phân loại giám sát.
    \item Cài đặt và kiểm chứng số học nghiêm ngặt toàn bộ pipeline bằng thư viện mảng NumPy thuần (đạt sai số tuyệt đối $< 10^{-11}$ và độ đồng thuận nhãn $100\%$ so với SciPy và scikit-learn), tính toán thống kê mô-men cục bộ bậc một và bậc hai thông qua Bảng tổng tích lũy với chi phí $\mathcal{O}(HW)$ độc lập với kích thước cửa sổ.
    \item Xây dựng không gian 8 đặc trưng vùng (4 đặc trưng hình học và 4 đặc trưng bức xạ/biên) giúp cắt giảm $76.6\%$ số vùng báo nhầm sinh ra sau phân ngưỡng (từ $14.89$ xuống $3.48$ FP/ảnh) trong khi giữ lại $97.8\%$ số đối tượng đúng của giai đoạn tạo ứng viên ($996/1\,018$ TP ứng viên).
    \item Thực hiện khảo sát thực nghiệm có kiểm soát trên lưới tham số quy mô lớn (110 cấu hình lọc/ngưỡng và 729 cấu hình mô hình vùng) trên tập kiểm định độc lập ($671$ ảnh Sentinel-1), chỉ rõ cơ chế suy biến của bộ lọc Lee về lọc trung bình hộp khi tham số cố định $L \in \{1, 4\}$ thấp hơn số looks hiệu dụng thực tế ($\text{ENL}$) của ảnh PNG đã qua tiền xử lý.
    \item Thẩm định khách quan theo giao thức ghép cặp đối tượng Hungarian ở ngưỡng $\IoU \ge 0.5$ và phân tầng theo 5 chế độ diện tích ($R_0 \to R_4$) trên tập kiểm thử độc lập $839$ ảnh Sentinel-1A và $776$ ảnh ALOS PALSAR, phân rã rõ hai nguồn lỗi trong chỉ số FP (báo nhầm thực sự trên nền biển sạch $\mathrm{Purity} \le 0.05$ và lỗi phân mảnh trên vùng dầu thật $\mathrm{Purity} \ge 0.5$) cũng như nút cổ chai độ thu hồi do bước phân ngưỡng toàn cục quyết định.
\end{itemize}
